\section{Исследование операций и задачи искусственного интеллекта. Экспертизы и неформальные процедуры. Автоматизация проектирования}

\subsection{Предмет и задачи исследования операций}

\textbf{Исследование операций} --- прикладная математическая дисциплина, изучающая методы обоснования и принятия оптимальных решений при управлении сложными системами.

Объектом исследования обычно является некоторая управляемая система, функционирование которой зависит от выбранного решения. Требуется определить такое решение, которое в заданных условиях обеспечивает наилучшее значение выбранного критерия.

В общем виде задачу можно представить как
\[
x^* = \arg\min_{x\in X} F(x),
\]
или
\[
x^* = \arg\max_{x\in X} F(x),
\]
где $x$ --- вектор управляемых параметров, $X$ --- множество допустимых решений, $F(x)$ --- критерий эффективности.

Типичная схема исследования операций включает:
\begin{enumerate}
	\item анализ реальной системы и постановку задачи;
	\item выделение управляемых и неуправляемых факторов;
	\item построение математической модели;
	\item выбор критерия эффективности;
	\item нахождение оптимального или рационального решения;
	\item анализ чувствительности решения;
	\item интерпретацию результатов и принятие решения.
\end{enumerate}

К основным классам задач исследования операций относятся:
\begin{itemize}
	\item распределение ограниченных ресурсов;
	\item транспортные задачи;
	\item задачи назначения;
	\item сетевые задачи и задачи о потоках;
	\item календарное и оперативное планирование;
	\item управление запасами;
	\item теория массового обслуживания;
	\item теория игр;
	\item управление проектами;
	\item имитационное моделирование.
\end{itemize}

Характерной особенностью исследования операций является системный подход: оптимизируется не отдельный элемент системы, а эффективность функционирования системы в целом.

\subsection{Модели принятия решений}

\textbf{Принятие решения} состоит в выборе одного варианта из множества допустимых альтернатив.

Формальная модель принятия решений содержит:
\begin{itemize}
	\item множество альтернатив $X$;
	\item множество возможных состояний внешней среды;
	\item критерии оценки решений;
	\item ограничения;
	\item информацию о последствиях выбора;
	\item правило выбора решения.
\end{itemize}

В простейшем случае имеется один критерий:
\[
F(x)\rightarrow \min,\qquad x\in X.
\]

При наличии нескольких критериев задача принимает вид
\[
F(x)=
\left(
F_1(x),F_2(x),\ldots,F_m(x)
\right),
\qquad x\in X.
\]

Одновременная оптимизация всех критериев обычно невозможна, поэтому применяют:
\begin{itemize}
	\item свёртку критериев:
	\[
	F(x)=\sum_{i=1}^{m}w_iF_i(x),
	\qquad
	w_i\geq 0,
	\qquad
	\sum_{i=1}^{m}w_i=1;
	\]
	\item выделение главного критерия и перевод остальных в ограничения;
	\item поиск множества Парето;
	\item методы экспертного ранжирования альтернатив.
\end{itemize}

Решение $x_1$ называется \textbf{Парето-доминирующим} решение $x_2$, если оно не хуже $x_2$ по всем критериям и хотя бы по одному критерию строго лучше.

В зависимости от имеющейся информации различают принятие решений:
\begin{itemize}
	\item \textbf{в условиях определённости} --- последствия каждого решения известны;
	\item \textbf{в условиях риска} --- возможные состояния и их вероятности известны;
	\item \textbf{в условиях неопределённости} --- вероятности состояний неизвестны;
	\item \textbf{в условиях конфликта} --- результат зависит также от решений других участников.
\end{itemize}

При наличии случайных факторов может оптимизироваться математическое ожидание:
\[
x^*=\arg\max_{x\in X}\mathsf{E}\,F(x,\xi),
\]
где $\xi$ --- случайный фактор.

\subsection{Оптимизационные задачи исследования операций}

Основной математический аппарат исследования операций связан с оптимизацией.

Общая задача математического программирования имеет вид
\[
F(x)\rightarrow\min,
\]
при ограничениях
\[
g_i(x)\leq 0,\qquad i=1,\ldots,m,
\]
\[
h_j(x)=0,\qquad j=1,\ldots,k.
\]

В зависимости от свойств модели возникают различные классы задач.

\textbf{Линейное программирование:}
\[
c^Tx\rightarrow\min,
\qquad
Ax\leq b.
\]

\textbf{Целочисленное программирование:}
\[
x_i\in\mathbb{Z}.
\]

Такие задачи возникают, например, когда переменная означает количество объектов или факт принятия решения:
\[
x_i\in\{0,1\}.
\]

\textbf{Нелинейное программирование} возникает, если целевая функция или ограничения нелинейны.

\textbf{Динамическое программирование} используется для многоэтапных процессов. Основная идея заключается в том, что оптимальное решение задачи строится из оптимальных решений её подзадач. Это выражается \textbf{принципом оптимальности Беллмана}.

\textbf{Сетевые модели} используются для задач поиска кратчайшего пути, максимального потока, минимального остовного дерева и других задач на графах.

Кроме аналитических и численных методов оптимизации широко применяется \textbf{имитационное моделирование}. Оно особенно полезно, когда явное аналитическое решение задачи получить невозможно.

\subsection{Основные задачи искусственного интеллекта}

\textbf{Искусственный интеллект} --- направление информатики, связанное с построением систем, способных решать задачи, для которых традиционно требуется интеллектуальная деятельность человека.

К основным задачам искусственного интеллекта относятся:
\begin{itemize}
	\item поиск решений в пространстве состояний;
	\item планирование действий;
	\item представление знаний;
	\item логический вывод;
	\item распознавание и классификация объектов;
	\item машинное обучение;
	\item обработка естественного языка;
	\item прогнозирование;
	\item поддержка принятия решений;
	\item управление интеллектуальными агентами и робототехническими системами.
\end{itemize}

Одним из классических подходов является представление задачи в виде пространства состояний.

Пусть $S$ --- множество возможных состояний системы, $s_0$ --- начальное состояние, $S_f$ --- множество целевых состояний. Необходимо найти последовательность допустимых переходов
\[
s_0\rightarrow s_1\rightarrow\ldots\rightarrow s_k,
\qquad
s_k\in S_f.
\]

Для поиска применяются:
\begin{itemize}
	\item поиск в ширину;
	\item поиск в глубину;
	\item поиск с ограничением глубины;
	\item эвристический поиск;
	\item алгоритм $A^*$.
\end{itemize}

В алгоритме $A^*$ вершина выбирается по оценке
\[
f(n)=g(n)+h(n),
\]
где $g(n)$ --- стоимость пути от начального состояния до $n$, а $h(n)$ --- эвристическая оценка стоимости пути от $n$ до цели.

Современные системы искусственного интеллекта часто используют машинное обучение. В этом случае зависимость
\[
y=f(x)
\]
не задаётся непосредственно программистом, а восстанавливается по обучающей выборке
\[
D=\{(x_i,y_i)\}_{i=1}^{N}.
\]

Исследование операций и искусственный интеллект имеют общую цель --- получение рационального решения. Однако в исследовании операций обычно предполагается наличие достаточно формализованной математической модели, тогда как методы ИИ применимы и к задачам, для которых знания представлены частично, эвристически или извлекаются из данных.

\subsection{Экспертные системы}

\textbf{Экспертная система} --- программная система, использующая знания специалистов некоторой предметной области для решения задач, требующих экспертных рассуждений.

Классическая экспертная система содержит:
\begin{itemize}
	\item \textbf{базу знаний};
	\item \textbf{рабочую память};
	\item \textbf{механизм логического вывода};
	\item подсистему приобретения и редактирования знаний;
	\item пользовательский интерфейс;
	\item подсистему объяснения полученного решения.
\end{itemize}

Знания могут представляться в виде:
\begin{itemize}
	\item продукционных правил;
	\item логических формул;
	\item фреймов;
	\item семантических сетей;
	\item онтологий.
\end{itemize}

Наиболее распространённая форма продукционного правила:
\[
\text{ЕСЛИ } A \text{, ТО } B.
\]

Например:
\[
\text{ЕСЛИ температура высокая И давление растёт,
	ТО состояние системы опасное.}
\]

Различают два основных способа применения правил.

\textbf{Прямой вывод} начинается с известных фактов. К ним последовательно применяются правила до получения новых фактов и заключения:
\[
\text{факты}\rightarrow\text{правила}\rightarrow\text{заключение}.
\]

\textbf{Обратный вывод} начинается с некоторой гипотезы. Система определяет, какие условия должны выполняться для её подтверждения, после чего пытается доказать эти условия.

Экспертные системы применяются для диагностики, классификации, интерпретации данных, выбора оборудования, планирования, проектирования и поддержки принятия решений.

Их важное достоинство --- возможность явно представить экспертные знания и объяснить логику получения результата.

Основной недостаток --- сложность извлечения знаний у экспертов и последующего сопровождения базы знаний. Это называют \textbf{проблемой приобретения знаний}.

\subsection{Экспертизы и неформальные процедуры}

Не все задачи принятия решений допускают полное математическое описание. Причинами могут быть:
\begin{itemize}
	\item недостаток количественных данных;
	\item наличие качественных критериев;
	\item уникальность рассматриваемой ситуации;
	\item высокая неопределённость;
	\item невозможность формализовать опыт специалиста.
\end{itemize}

В таких случаях используются \textbf{экспертные оценки} и \textbf{неформальные процедуры}.

Экспертная оценка заключается в получении и обработке суждений специалистов предметной области.

Основными методами являются:
\begin{itemize}
	\item ранжирование альтернатив;
	\item непосредственная балльная оценка;
	\item попарное сравнение;
	\item анкетирование;
	\item интервьюирование;
	\item метод Делфи;
	\item коллективное обсуждение;
	\item построение сценариев.
\end{itemize}

При ранжировании эксперт располагает объекты в порядке предпочтения:
\[
x_{i_1}\succ x_{i_2}\succ\ldots\succ x_{i_n}.
\]

При использовании нескольких экспертов их оценки могут агрегироваться, например,
\[
R_i=\sum_{j=1}^{m}w_jr_{ij},
\]
где $r_{ij}$ --- оценка объекта $i$ экспертом $j$, $w_j$ --- вес компетентности эксперта.

Для попарных сравнений может использоваться матрица
\[
A=(a_{ij}),
\]
где $a_{ij}$ характеризует предпочтительность альтернативы $i$ относительно альтернативы $j$.

\textbf{Метод Делфи} основан на последовательном анонимном опросе группы экспертов. После каждого раунда экспертам предоставляется обобщённая информация о результатах предыдущего раунда, что позволяет постепенно сближать оценки.

Неформальные процедуры не означают произвольное принятие решений. Их задача --- формализовать настолько, насколько возможно, знания и опыт специалистов там, где построение полной математической модели затруднено.

На практике эффективным является сочетание:
\[
\text{математическая модель}
+
\text{данные}
+
\text{экспертные знания}.
\]

\subsection{Системы автоматизированного проектирования}

\textbf{Система автоматизированного проектирования (САПР)} --- организационно-техническая система, предназначенная для автоматизации процессов проектирования объектов и систем.

В англоязычной терминологии часто используется обозначение \textbf{CAD} (\textit{Computer-Aided Design}).

Процесс проектирования можно представить как поиск параметров объекта
\[
x=(x_1,\ldots,x_n)
\]
при выполнении совокупности требований
\[
g_i(x)\leq 0
\]
и достижении требуемых характеристик
\[
F(x)\rightarrow\min
\quad\text{или}\quad
F(x)\rightarrow\max.
\]

САПР могут обеспечивать:
\begin{itemize}
	\item геометрическое моделирование;
	\item параметрическое проектирование;
	\item ведение проектной документации;
	\item инженерные расчёты;
	\item моделирование поведения проектируемого объекта;
	\item проверку ограничений;
	\item оптимизацию параметров конструкции;
	\item хранение и повторное использование проектных данных.
\end{itemize}

В инженерной практике обычно различают:
\begin{itemize}
	\item \textbf{CAD} --- построение и редактирование модели объекта;
	\item \textbf{CAE} (\textit{Computer-Aided Engineering}) --- инженерный анализ и вычислительное моделирование;
	\item \textbf{CAM} (\textit{Computer-Aided Manufacturing}) --- автоматизированная подготовка производства.
\end{itemize}

При параметрическом проектировании модель задаётся набором параметров. Изменение параметров автоматически приводит к перестроению геометрии и связанных характеристик объекта.

Это позволяет автоматически исследовать множество вариантов проекта и решать задачи оптимального проектирования.

\subsection{Интеграция оптимизации, искусственного интеллекта и проектирования}

Современная автоматизация проектирования объединяет методы математического моделирования, оптимизации, искусственного интеллекта и экспертного анализа.

Типичный автоматизированный цикл имеет вид:
\[
\boxed{
	\text{параметры проекта}
	\rightarrow
	\text{математическая модель}
	\rightarrow
	\text{расчёт}
	\rightarrow
	\text{оценка критериев}
	\rightarrow
	\text{оптимизация}
	\rightarrow
	\text{новые параметры}
}
\]

Цикл повторяется до достижения заданного критерия остановки.

Например, при проектировании конструкции требуется подобрать вектор параметров $x$ так, чтобы
\[
m(x)\rightarrow\min,
\]
при ограничениях
\[
\sigma_{\max}(x)\leq[\sigma],
\qquad
u_{\max}(x)\leq[u],
\]
где $m$ --- масса конструкции, $\sigma_{\max}$ --- максимальное напряжение, $u_{\max}$ --- максимальное перемещение.

Математическая оптимизация позволяет автоматически искать лучшие значения параметров.

Методы искусственного интеллекта могут использоваться для:
\begin{itemize}
	\item выбора первоначальной конфигурации;
	\item прогнозирования характеристик проектируемого объекта;
	\item поиска закономерностей в результатах расчётов;
	\item построения приближённых моделей;
	\item генерации вариантов проекта;
	\item обнаружения ошибочных или недопустимых решений;
	\item использования накопленного опыта предыдущих проектов.
\end{itemize}

Экспертные знания позволяют учитывать требования, которые трудно представить в виде точных математических ограничений.

Таким образом, современная интеллектуальная САПР может рассматриваться как система поддержки принятия проектных решений, в которой объединяются
\[
\boxed{
	\text{модель}
	+
	\text{оптимизация}
	+
	\text{данные}
	+
	\text{экспертные знания}
	+
	\text{ИИ}.
}
\]

При этом окончательное решение во многих задачах остаётся за человеком. Автоматизированная система перебирает и анализирует большое количество вариантов, проверяет ограничения и предлагает рациональные решения, а специалист оценивает факторы, которые невозможно полностью формализовать.

\newpage